\section{Parallel Universes and Eternal Return}

\textbf{To the May Bride}

\begin{quotex}
I know that I will meet you again and that everything will happen once again exactly as it did so long ago. Except that this time I will not allow you to die. I will hold you in my arms, defending you against the dark waters of death. Because this time I will remember everything. \emph{I will remember that you have already died.} But ... will I remember?

\flright{\textsc{Miguel Serrano}, \emph{Nos: Book of the Resurrection}}

Because You are unlimited, neither the lords of heaven nor even You Yourself can ever reach the end of Your glories. The \emph{countless universes}, each enveloped in its shell, are compelled by the wheel of time to wander within You, like particles of dust blowing about in the sky. The srutis, following their method of eliminating everything separate from the Supreme, become successful by revealing You as their final conclusion 
\flright{\textit{Bhagavata Purana}, 10.87.41}
\end{quotex}


Quantum Mechanics often mimics traditional metaphysics, as we have see the example of the ultimate indeterminateness or formless of matter. Such examples do not confirm metaphysics, but rather are what would be expected given the truth of metaphysical principles. Another such idea is ``parallel universes'' or the multiverse. In this conception, the universe we find ourselves in is not the only one. Every time a choice is made, another universe comes into existence where the opposite choice had been made.

However, the continual creation of new universes in time is not precisely the traditional teaching. First of all, the QM version does not take into account the notion of compossibilities, which limits to some degree seemingly possible universes. The other is that other universes may be quite different from the one we are conscious of. Although the idea of possible universes is not encountered often in religious thinking, it does have a pedigree. The Islamic philosopher Fakhr al-Din al-Razi\footnote{\url{http://en.wikipedia.org/wiki/Fakhr\_al-Din\_al-Razi}} came to this conclusion based on verses in the Koran regarding Allah as Lord of the Words plural . He wrote:

\begin{quotex}
It is established by evidence that there exists beyond the world a void without a terminal limit (khala' la nihayata laha), and it is established as well by evidence that God Most High has power over all contingent beings (al-mumkinat). Therefore He the Most High has the power (qadir) to create a thousand thousand worlds (alfa alfi `awalim) beyond this world such that each one of those worlds be bigger and more massive than this world as well as having the like of what this world has of the throne (al-arsh), the chair (al-kursiyy), the heavens (al-samawat) and the earth (al-ard), and the sun (al-shams) and the moon (al-qamar). The arguments of the philosophers (dala'il al-falasifah) for establishing that the world is one are weak, flimsy arguments founded upon feeble premises.
\end{quotex}


The Hermetist Giordano Bruno\footnote{\url{http://en.wikipedia.org/wiki/Giordano_Bruno}} makes an argument for the existence of multiple worlds based on the infinity of God. In \emph{Ash Wednesday Supper} he writes:

\begin{quotex} He made his affirmation that the universe is infinite; that it consists of an immense ethereal region; that it is like a vast sky of space in whose bosom are the heavenly bodies... that the moon, the sun, and innumerable other bodies are in this ethereal region, and the earth also; that is is not to be believed that there is any firmament, base, or foundation to which are fixed these great animals which form the constitution of the universe the infinite material of the infinite divine potency.
\end{quotex}

Friedrich Nietzsche claims the same world repeats:

\begin{quotex}
What, if some day or night a demon were to steal after you into your loneliest loneliness and say to you: `This life as you now live it and have lived it, you will have to live once more an innumerable times more' ... Would you not throw yourself down and gnash your teeth and curse the demon who spoke thus? Or have you once experienced a tremendous moment when you would have answered him: `You are a god and never have I heard anything more divine.'
\flright{\textit{The Gay Science}, 341}
\end{quotex}


The concept of Eternal Return is illustrated in the diagram. A life is represented by the circle, and the world by the plane. Life begins, it completes, and at some future time the life will begin again at the same point, and repeat the same circle. It assumes that space is finite and time infinite; thus every spatial configuration will eventually have to be repeated. Note that the center of the circle is on the material plane, so there is no sense of transcendence. Also, as Guenon points out in \emph{The Symbolism of the Cross}, Infinity is incompatible with the idea of repetition. Every possibility of manifestation happens once, and once only.

The second diagram is adapted from Rene Guenon. Each plane represents a world and a state of being. Due to the infinitude of possibilities, no state is repeated. Thus a being that manifests in the human state will never repeat that state. This obviously precludes reincarnation in the sense of a being repeating earthbound human lives. Again, we take the circle to represent a life along with its center. For reincarnation to occur, there would have to be a new circle with a new center; this is hard to conceive.

Unlike eternal return, therefore, Guenon proposes as the Traditional view (see the above quotes) the idea of multiple worlds, represented by the various planes. As a circle closes, instead of closing on itself --- and ending its possibilities --- it connects to the beginning of a circle in the world above it. This continues, forming a helix with a common center. Since birth and death are the end points of the circle, we see that death in one world corresponds to birth in another, and vice versa. Now, this death is not necessarily a physical death, but may include post-mortem states appropriate to that plane of existence. Nevertheless, the the possibilities of that world have been played out, even those states will terminate.

Guenon does not, and cannot, describe the features of another world, but we may assume that a plane that is infinitesimally close to ours, may be similar in many ways but also slightly different. In this sense, we may regard it as a form of reincarnation in that a being is incarnated in a material world, though not the same as the previous one. (Note that temporal sequence does not have the precise meaning in the conventional way, but it is easier to describe it in those terms). Understood in this way, this doctrine is similar to that of Fourth Way teachings of Ouspensky.

\textit{Many thanks to Prasad Sarangapani for creating the geometric diagrams.}


\flrightit{Posted on 2011-05-24 by Cologero}

\begin{center}* * *\end{center}

\begin{footnotesize}
\begin{sffamily}

\texttt{Izak on 2011-05-28 at 17:47 said:}

The second diagram is somewhat incomplete, it seems. In Symbolism of the Cross (which, at the moment I admit I don't have handy), I recall Guenon using the image of an upward spiral contained within a sphere that possesses a three-dimensional (that is, six-pointed) cross. The spiral itself would seem to be planar, not linear. I envisioned this image, basically.\footnote{\url{http://lh4.ggpht.com/\_mlEZC-GqQwc/THgRr\_aw--I/AAAAAAAAAY8/k\_UPz-XU8-c/spiral\_curve\%5B1\%5D.png}}

but with planar coordinates filled in from the vertical pole (which is impossible, I think, to actually draw). That way, it allows us to recognize the indefinitude of horizontal planes, because there is no precise number of spiral swirls and thus no points of degree separating each spiral swirl from the other (one recalls Zeno's argument about infinite space existing between two points). My point is not to say that the diagram is wrong, but that it does not seem to be the complete picture of what Guenon was describing, which is probably impossible to actually convey through illustration anyhow. The diagram here seems like more of a step in a process of mentally constructing it. It is frustrating that Guenon himself doesn't seem to have included his own diagrams.

Also he takes pains to assure the reader that he is not talking about an ``infinitude'' of possibilities but an indefinite number. He seems very concerned about avoiding discussions of infinite space and states of being.

Anyhow, good discussion on metaphysics.

\hfill

\texttt{Cologero on 2011-05-28 at 18:44 said:}

Thanks for the diagram, it takes some mental gymnastics to visualize the diagrams from Guenon's verbose description. I did, however, point out that the diagram provided was an adaptation to demonstrate the point. Oddly enough, in traditional teachings there would be nothing \emph{but} diagrams, avoiding the very necessity of the text. This would go along with oral teachings related to the diagrams. It is one of our goals to realize all the diagrams described in Symbolism of the Cross and we are always looking for help.

Guenon defines ``infinite'' very narrowly to mean unbounded in the metaphysical sense so he never uses that word as a mathematician might. By ``indefinite'' he means ``uncountable'', or ``uncountably infinite'' (the last word used in the mathematical sense.) So Universal Possibility is unlimited, or infinite (see Multiple States of Being), although the specific possibilities for a being are not (the True Man actualizes all his possibilities). Nevertheless, the possible states of being are uncountable.

\hfill

\end{sffamily}
\end{footnotesize}

